%========================================================================
\section{Conclusion}
% TIMING ~3 min  |  Refermer l'axe. PAS de CDI. Finir sur la substance et la trajectoire.
%========================================================================

% --- C1 : Schema de l'evolution QA -> Produit -------------------------
\begin{frame}{Une évolution des responsabilités}
  % VOIX ~1 min 30. Le point d'intervention remonte la chaine.
  \begin{center}
  \begin{tikzpicture}[
    etape/.style={rectangle, rounded corners, draw=ClaudeDark, very thick,
                  fill=ClaudeCream, text width=2.7cm, align=center, minimum height=1.7cm},
    fleche/.style={-{Stealth[length=3mm]}, very thick, ClaudeOrange}
  ]
    \node[etape] (a) {\textbf{Vérifier}\\[2pt]\scriptsize QA : Homologation, TNR, documentation};
    \node[etape, right=0.7cm of a] (b) {\textbf{Définir}\\[2pt]\scriptsize Product : rédaction de tickets};
    \node[etape, right=0.7cm of b] (c) {\textbf{Orienter}\\[2pt]\scriptsize Discovery : benchmark, audits};
    \draw[fleche] (a) -- (b);
    \draw[fleche] (b) -- (c);
  \end{tikzpicture}
  \end{center}
  \vskip 6pt
  \begin{block}{Le point commun}
    \footnotesize Transformer un \textbf{processus non structuré} en quelque chose de \textbf{reproductible et transmissible} en \textbf{utilisant l'IA} tout en gardant la main sur la décision finale.
  \end{block}
\end{frame}

% --- C2 : Avenir + remerciements --------------------------------------
\begin{frame}{Et maintenant}
  % VOIX ~1 min 30. Trajectoire + remerciements. PAS de CDI.
  \begin{itemize}\footnotesize
    \item \textbf{Ce que le stage a confirmé} : une trajectoire vers la gestion de projet et le Product
    \item \textbf{Poursuivre} : des missions de Discovery et d'intégration de l'IA dans le quotidien
  \end{itemize}
  \vskip 8pt
  \begin{exampleblock}{Remerciements}
    \footnotesize Merci à \textbf{Alexandre Séjourné} pour son encadrement et à toute l'équipe d'Ecocea La Fabrique pour leur confiance.
  \end{exampleblock}
\end{frame}
